\section{\ENG{xAPI} / \ENG{Tin} \ENG{Can} \ENG{API} (2011-2016)}
\label{sec:xapi_tincan}

Το \ENG{Experience} \ENG{API} (\ENG{xAPI}), γνωστό και ως \ENG{Tin} \ENG{Can} \ENG{API}, αποτελεί θεμελιώδη εξέλιξη στα \ENG{e-learning} \ENG{standards}, σηματοδοτώντας τη μετάβαση από \ENG{monolithic}, \ENG{browser-centric} \ENG{architectures} σε \ENG{distributed}, \ENG{cloud-ready} \ENG{systems} που μπορούν να καταγράφουν μαθησιακές εμπειρίες σε οποιοδήποτε περιβάλλον.
\subsection{Γέννηση και Ανάπτυξη του \ENG{xAPI}}
\label{subsec:xapi_genesis}

\subsubsection{\ENG{Project} \ENG{Tin} \ENG{Can} (2011)}

Το 2011, το \ENG{Advanced} \ENG{Distributed} \ENG{Learning} (\ENG{ADL}) \ENG{Initiative} εξέδωσε ένα \textit{\ENG{Broad} \ENG{Agency} \ENG{Announcement}} (\ENG{BAA}) αναζητώντας λύσεις για τους περιορισμούς του \ENG{SCORM.} Το \ENG{ADL} αναγνώριζε ότι το \ENG{SCORM}, σχεδιασμένο για \ENG{web} \ENG{browsers} και \ENG{LMS-centric} περιβάλλοντα, δεν μπορούσε να ανταποκριθεί στις σύγχρονες μαθησιακές ανάγκες:

Περιορισμοί \ENG{SCORM} που οδήγησαν στο \ENG{xAPI}: η εξάρτηση από \ENG{web} \ENG{browser} για την εκτέλεση του \ENG{content}, ο εγκλωβισμός των δεδομένων στο \ENG{LMS}, το περιορισμένο \ENG{data} \ENG{model} με σταθερά πεδία που δεν επεκτείνονται, η απουσία \ENG{mobile} \ENG{support} για \ENG{native} \ENG{mobile} \ENG{apps}, η έλλειψη \ENG{offline} \ENG{tracking} χωρίς συνεχή σύνδεση με \ENG{LMS} και ο περιορισμός σε ένα μόνο \ENG{learning} \ENG{context} χωρίς συνδυασμό πολλαπλών περιβαλλόντων.

Η \ENG{Rustici} \ENG{Software} ανέπτυξε το \textit{\ENG{Project} \ENG{Tin} \ENG{Can}}, το οποίο μετονομάστηκε σε \ENG{Experience} \ENG{API} (\ENG{xAPI}). Το όνομα "\ENG{Tin} \ENG{Can}" προέρχεται από την παιδική ιδέα της επικοινωνίας με "\ENG{tin} \ENG{can} \ENG{telephone}" - απλή, \ENG{distributed} \ENG{communication.}

\subsubsection{Χρονολόγιο Εξέλιξης}

Το χρονολόγιο εξέλιξης περιλαμβάνει την κυκλοφορία της \ENG{xAPI} 1.0.0 \ENG{specification} τον Απρίλιο του 2013, την ταχεία υιοθέτηση από \ENG{vendors} και \ENG{LRS} \ENG{providers} το 2013-2014, την έκδοση \ENG{xAPI} 1.0.2 το 2015 (\ENG{bug} \ENG{fixes}, \ENG{clarifications}), την \ENG{xAPI} 1.0.3 το 2016 (\ENG{minor} \ENG{improvements}) και, το 2023, την επίσημη αναγνώριση ως \ENG{IEEE} 9274.1.1-2023 \ENG{standard}.

\subsection{Αρχιτεκτονική \ENG{xAPI}}
\label{subsec:xapi_architecture}

\subsubsection{Βασική Ιδέα: \ENG{Actor-Verb-Object} \ENG{Statements}}

Το \ENG{xAPI} βασίζεται σε απλή ιδέα: καταγραφή εμπειριών με \ENG{statements} της μορφής:

\begin{center}
\textbf{"\ENG{Actor} \ENG{Verb} \ENG{Object}"}
\end{center}

Παραδείγματα: "\ENG{John} \textit{\ENG{completed}} \ENG{Module} 1", "\ENG{Mary} \textit{\ENG{watched}} \ENG{Video} \ENG{Tutorial}", "\ENG{Team} \ENG{A} \textit{\ENG{achieved}} \ENG{Sales} \ENG{Goal}" και "\ENG{Student} \textit{\ENG{failed}} \ENG{Final} \ENG{Exam}".

\subsubsection{Ανατομία \ENG{xAPI} \ENG{Statement}}

Ένα \ENG{xAPI} \ENG{statement} είναι \ENG{JSON} \ENG{object} με τα εξής κύρια στοιχεία:

{\selectlanguage{english}\fontencoding{T1}\selectfont
\begin{verbatim}
{
  "actor": {
    "objectType": "Agent",
    "name": "John Doe",
    "mbox": "mailto:john.doe@example.com"
  },
  "verb": {
    "id": "http://adlnet.gov/expapi/verbs/completed",
    "display": {"en-US": "completed"}
  },
  "object": {
    "objectType": "Activity",
    "id": "http://example.com/activities/module1",
    "definition": {
      "name": {"en-US": "Introduction to SCORM"},
      "description": {"en-US": "Basic SCORM concepts"}
    }
  },
  "result": {
    "score": {"scaled": 0.85, "raw": 85, "min": 0, "max": 100},
    "completion": true,
    "success": true,
    "duration": "PT2H30M"
  },
  "context": {
    "registration": "a4c5e0b3-...",
    "contextActivities": {
      "parent": [{"id": "http://example.com/courses/scorm101"}]
    },
    "platform": "Example LMS",
    "language": "en-US"
  },
  "timestamp": "2026-01-01T10:30:00Z",
  "stored": "2026-01-01T10:30:01Z"
}
\end{verbatim}
}

\begin{table}[H]
\centering
\caption{Βασικά στοιχεία ενός \ENG{xAPI} \ENG{statement}}
\label{tab:xapi_statement_elements}
\small
\setlength{\tabcolsep}{5pt}
\renewcommand{\arraystretch}{1.25}
\begin{chaptertablebox}
\begin{tabular}{|c|p{3.7cm}|p{8.2cm}|}
\hline
\textbf{Α/Α} & \textbf{Στοιχείο} & \textbf{Περιγραφή} \\
\hline
\hline
1 & \textbf{\ENG{Actor}} (Υποχρεωτικό) & Ο μαθητής ή η οντότητα που εκτελεί τη δράση. Μπορεί να είναι \ENG{Agent} (μονό άτομο με \ENG{mbox}, \ENG{mbox}\_\ENG{sha1sum}, \ENG{openid}, \ENG{account}) ή \ENG{Group} (ομάδα ατόμων, \ENG{identified} ή \ENG{anonymous}). \\
\hline
2 & \textbf{\ENG{Verb}} (Υποχρεωτικό) & Η δράση που εκτελείται. Προσδιορίζεται από \ENG{id} (\ENG{URI} που ορίζει το \ENG{verb}) και \ENG{display} (\ENG{human-readable} σε πολλές γλώσσες), με ενδεικτικά παραδείγματα \ENG{completed}, \ENG{passed}, \ENG{failed}, \ENG{answered}, \ENG{attempted}, \ENG{experienced}, \ENG{watched}, \ENG{listened}. \\
\hline
3 & \textbf{\ENG{Object}} (Υποχρεωτικό) & Η δραστηριότητα, \ENG{agent} ή \ENG{statement} που αποτελεί το αντικείμενο. Μπορεί να είναι \ENG{Activity} (εκπαιδευτική δραστηριότητα), \ENG{Agent/Group} (άλλος χρήστης), \ENG{StatementRef} (αναφορά σε άλλο \ENG{statement}) ή \ENG{SubStatement} (εμφωλευμένο \ENG{statement}). \\
\hline
4 & \textbf{\ENG{Result}} (Προαιρετικό) & Αποτέλεσμα της δράσης. Περιλαμβάνει \ENG{score} (βαθμολογία \ENG{scaled}, \ENG{raw}, \ENG{min}, \ENG{max}), \ENG{completion} (\ENG{true/false}), \ENG{success} (\ENG{true/false}), \ENG{duration} σε \ENG{ISO} 8601 \ENG{format} (\ENG{PT2H30M} = 2 ώρες 30 λεπτά), \ENG{response} και \ENG{extensions} (\ENG{custom} \ENG{data}). \\
\hline
5 & \textbf{\ENG{Context}} (Προαιρετικό) & Πλαίσιο της εμπειρίας. Περιλαμβάνει \ENG{registration} (\ENG{UUID} για \ENG{grouping} \ENG{statements}), \ENG{contextActivities} (\ENG{parent}, \ENG{grouping}, \ENG{category}, \ENG{other}), \ENG{instructor}, \ENG{team}, \ENG{platform}, \ENG{language}, \ENG{statement} (\ENG{reference} σε σχετικό \ENG{statement}) και \ENG{extensions} (\ENG{custom} \ENG{context} \ENG{data}). \\
\hline
6 & \textbf{\ENG{Timestamp}} (Προαιρετικό) & Δηλώνει πότε συνέβη η εμπειρία, σε μορφή \ENG{ISO} 8601. \\
\hline
7 & \textbf{\ENG{Stored}} (Αυτόματο από \ENG{LRS}) & Δηλώνει πότε αποθηκεύτηκε το \ENG{statement}. \\
\hline
8 & \textbf{\ENG{Authority}} (Αυτόματο) & Δηλώνει ποιος έστειλε το \ENG{statement}. \\
\hline
\end{tabular}
\end{chaptertablebox}
\end{table}


\subsection{\ENG{Learning} \ENG{Record} \ENG{Store} (\ENG{LRS})}
\label{subsec:lrs}

Το \ENG{Learning} \ENG{Record} \ENG{Store} (\ENG{LRS}) είναι κεντρικό στοιχείο της αρχιτεκτονικής \ENG{xAPI.} Αποτελεί το σύστημα αποθήκευσης \ENG{xAPI} \ENG{statements}~\cite{xapicom2024lrs,rustici2023lrs}.

\subsubsection{Χαρακτηριστικά \ENG{LRS}}

Τα βασικά χαρακτηριστικά και οι κύριες κατηγορίες \ENG{LRS} συνοψίζονται στον Πίνακα~\ref{tab:lrs_characteristics}.

\begin{table}[H]
\centering
\caption{Χαρακτηριστικά και βασικές κατηγορίες \ENG{LRS}}
\label{tab:lrs_characteristics}
\small
\setlength{\tabcolsep}{5pt}
\renewcommand{\arraystretch}{1.25}
\begin{chaptertablebox}
\begin{tabular}{|p{4cm}|p{11cm}|}
\hline
\textbf{Χαρακτηριστικό} & \textbf{Περιγραφή / Παραδείγματα} \\
\hline
\ENG{RESTful API} & Παρέχει \ENG{HTTP endpoints} όπως:  \ENG{\texttt{POST /statements}, \texttt{GET /statements}, \texttt{PUT /statements}, \texttt{GET /activities/state}, \texttt{GET /activities/profile}, \texttt{GET /agents/profile}}. \\
\hline
\ENG{Authentication} & Υποστηρίζει \ENG{OAuth 2.0} (\ENG{preferred}), \ENG{HTTP Basic Authentication}, \ENG{Certificate-based authentication}. \\
\hline
\ENG{Standalone LRS} & Ανεξάρτητο σύστημα (\ENG{cloud} ή \ENG{on-premise}), π.χ. \ENG{Learning Locker}, \ENG{Watershed}, \ENG{Veracity Learning}, \ENG{SCORM Cloud}. Προσφέρει πλήρη έλεγχο και \ENG{advanced analytics}. \\
\hline
\ENG{Integrated LRS} & Ενσωματωμένο σε \ENG{LMS}, π.χ. \ENG{Moodle xAPI plugin}, \ENG{Totara LMS}. Προσφέρει απλοποιημένη ενσωμάτωση με περιορισμένες δυνατότητες. \\
\hline
\ENG{Cloud LRS Services} & \ENG{SaaS offerings}, π.χ. \ENG{Rustici SCORM Cloud}, \ENG{Yet Analytics Watershed}. Παρέχει \ENG{low maintenance}, \ENG{scalability} και \ENG{analytics dashboards}. \\
\hline
\end{tabular}
\end{chaptertablebox}
\end{table}


\subsection{Τα "4 \ENG{Freedoms}" του \ENG{xAPI}}
\label{subsec:xapi_freedoms}

Το \ENG{xAPI} σχεδιάστηκε γύρω από τέσσερις θεμελιώδεις ελευθερίες~\cite{adl2013xapi_spec}:

\begin{itemize}
  \item \ENG{Freedom of Statement:} Επιτρέπει την καταγραφή οποιασδήποτε εμπειρίας, όχι μόνο \ENG{course completions}.
  \item \ENG{Freedom of History:} Τα δεδομένα αποθηκεύονται μόνιμα στο \ENG{LRS}, και δεν χάνονται όπως τα \ENG{session data}.
  \item \ENG{Freedom of Device:} Εξασφαλίζει ότι το \ENG{xAPI} λειτουργεί σε \ENG{desktop}, \ENG{mobile}, \ENG{tablets}, \ENG{IoT} και \ENG{wearables}.
  \item \ENG{Freedom of Workflow:} Δεν περιορίζεται σε \ENG{linear workflows}, υποστηρίζοντας \ENG{blended}, \ENG{social} και \ENG{informal learning}.
\end{itemize}

\subsection{\ENG{xAPI} \ENG{Profiles} και \ENG{Recipes}}
\label{subsec:xapi_profiles}

Για να εξασφαλιστεί \ENG{consistency}, το \ENG{xAPI} χρησιμοποιεί \ENG{Profiles}, δηλαδή έγγραφα που ορίζουν συγκεκριμένα \ENG{verbs} για το \ENG{domain}, \ENG{activity} \ENG{types}, \ENG{statement} \ENG{templates} και \ENG{rules} ή \ENG{patterns}. Ενδεικτικά, το \ENG{Video} \ENG{Profile} ορίζει \ENG{verbs} για \ENG{played}, \ENG{paused}, \ENG{seeked} και \ENG{volume-change}, το \ENG{Serious} \ENG{Games} \ENG{Profile} καλύπτει \ENG{game-specific} \ENG{interactions}, ενώ το \ENG{cmi5} \ENG{Profile} παρέχει \ENG{SCORM-like} \ENG{structure} με \ENG{xAPI}.

\subsection{Ασφάλεια και Προστασία Δεδομένων}
\label{subsec:xapi_security}

Το \ENG{xAPI} απαιτεί \ENG{HTTPS} για \ENG{production} και \ENG{TLS} 1.2+ \ENG{encryption}. Υποστηρίζει \ENG{OAuth} 2.0 με \ENG{scopes} (\ENG{read}, \ENG{write}, \ENG{define}), \ENG{HTTP} \ENG{Basic} με \ENG{SSL} και \ENG{certificate-based} \ENG{authentication}. Η \ENG{privacy} περιλαμβάνει \ENG{granular} \ENG{permissions} (\ENG{read/write/define}), δυνατότητες \ENG{anonymization}, συμμόρφωση με \ENG{GDPR} και υποστήριξη για \ENG{right} \ENG{to} \ENG{erasure}.

\subsection{Υιοθέτηση και \ENG{Ecosystem} (2016-2026)}
\label{subsec:xapi_adoption}

Η υιοθέτηση περιλαμβάνει χρήση από \ENG{Fortune} 500 εταιρείες, \ENG{government/military} \ENG{applications}, \ENG{healthcare} \ENG{training} \ENG{programs} και \ENG{higher} \ENG{education} \ENG{initiatives}. Ενδεικτικά εργαλεία συγγραφής περιλαμβάνουν τα \ENG{Articulate} \ENG{Storyline} 360 (\ENG{xAPI} \ENG{support}), \ENG{Adobe} \ENG{Captivate}, \ENG{Adapt} \ENG{Learning} \ENG{Framework} και \ENG{Lectora}. Οι ενδεικτικοί \ENG{LRS} \ENG{providers} περιλαμβάνουν τους \ENG{Learning} \ENG{Locker} (\ENG{open} \ENG{source}), \ENG{Watershed} \ENG{by} \ENG{Rustici}, \ENG{Veracity} \ENG{Learning} και \ENG{SCORM} \ENG{Cloud}. Το \ENG{xAPI} έχει γίνει το \ENG{standard} επιλογής για οργανισμούς που χρειάζονται \ENG{comprehensive} \ENG{learning} \ENG{analytics} και \ENG{tracking} πέρα από τα \ENG{traditional} \ENG{e-learning} \ENG{courses}.

% \ENG{Figure}: \ENG{xAPI} \ENG{Statement} \ENG{Anatomy}
\begin{figure}[H]
  \centering
  \begin{chapterfigurebox}
\begin{tikzpicture}[
  >=Latex,
  every node/.style={font=\small},
  corebox/.style={
    rectangle,
    draw,
    thick,
    rounded corners,
    minimum width=4cm,
    minimum height=2.8cm,
    text width=3.8cm,
    align=left,
    inner sep=7pt,
    font=\ttfamily\small
  },
  metabox/.style={
    rectangle,
    draw,
    thick,
    rounded corners,
    minimum width=13cm,
    minimum height=2.8cm,
    text width=12.5cm,
    align=left,
    inner sep=7pt,
    font=\ttfamily\small
  },
  boxlabel/.style={font=\bfseries\small, align=center},
  arrowlabel/.style={font=\footnotesize, fill=white, inner sep=1pt},
  note/.style={font=\small, align=center, text width=12.6cm}
]

% Title
\node[font=\large\bfseries, text=xapiorange] (title) at (0,4.5)
  {\ENG{xAPI} \ENG{Statement} \ENG{Structure}};

% Actor
\node[corebox, fill=blue!12] (actor) at (-6,1.5) {%
"\ENG{actor}": \{\\
\quad "\ENG{type}": "\ENG{Agent}",\\
\quad "\ENG{name}": "\ENG{John Doe}",\\
\quad "\ENG{mbox}": "\ENG{mailto:john@example.com}"\\
\}
};
\node[boxlabel, text=blue!70!black] at ($(actor.north)+(0,0.45)$)
  {\ENG{ACTOR} (Ποιος)};

% Verb
\node[corebox, fill=xapiorange!20] (verb) at (0,1.5) {%
"\ENG{verb}": \{\\
\quad "\ENG{id}": "\ENG{adlnet.gov/.../completed}",\\
\quad "\ENG{display}": \{"\ENG{en-US}": "\ENG{completed}"\}\\
\}
};
\node[boxlabel, text=xapiorange!80!black] at ($(verb.north)+(0,0.45)$)
  {\ENG{VERB} (Τι έκανε)};

% Object
\node[corebox, fill=green!12] (object) at (6,1.5) {%
"\ENG{object}": \{\\
\quad "\ENG{type}": "\ENG{Activity}",\\
\quad "\ENG{id}": "\ENG{example.com/course-101}",\\
\quad "\ENG{name}": "\ENG{Course 101}"\\
\}
};
\node[boxlabel, text=green!50!black] at ($(object.north)+(0,0.45)$)
  {\ENG{OBJECT} (Σε τι)};

% Arrows
\draw[->, thick, gray!70] (actor.east) -- node[arrowlabel, above] {\ENG{does}} (verb.west);
\draw[->, thick, gray!70] (verb.east) -- node[arrowlabel, above] {\ENG{to}} (object.west);

% Optional/automatic fields
\node[metabox, fill=gray!10] (meta) at (0,-2.8) {%
"\ENG{result}": \{"\ENG{score}": 0.85, "\ENG{success}": \ENG{true}, "\ENG{completion}": \ENG{true}\}\\
"\ENG{context}": \{"\ENG{platform}": "\ENG{Example LMS}", "\ENG{instructor}": "...", "\ENG{registration}": "\ENG{uuid}"\}\\
"\ENG{timestamp}": "\ENG{2024-01-15T10:30:00Z}",\ "\ENG{stored}": "\ENG{2024-01-15T10:30:05Z}",\ "\ENG{authority}": "\ENG{LRS/client}"
};
\node[boxlabel, text=black!70] at ($(meta.north)+(0,0.45)$)
  {\ENG{OPTIONAL / AUTOMATIC FIELDS}};

% Footnote
\node[note] (footnote) at (0,-5.3) {%
Το ελάχιστο υποχρεωτικό \ENG{xAPI} \ENG{statement} αποτελείται από \ENG{Actor}, \ENG{Verb} και \ENG{Object}, ενώ τα πεδία \ENG{Result}, \ENG{Context}, \ENG{Timestamp}, \ENG{Stored} και \ENG{Authority} προσθέτουν επιπλέον πληροφορία καταγραφής.
};

% Background box
\begin{pgfonlayer}{background}
  \node[
    draw=xapiorange,
    very thick,
    rounded corners,
    fill=xapiorange!4,
    fit=(title)(actor)(verb)(object)(meta)(footnote),
    inner sep=12pt
  ] {};
\end{pgfonlayer}

\end{tikzpicture}
\end{chapterfigurebox}

  \caption{Ανατομία \ENG{xAPI} \ENG{Statement} - \ENG{Actor-Verb-Object} \ENG{Structure}}
  \label{fig:xapi_statement}
\end{figure}

